\bookchapter{第三十六章：スマホとスマートホーム——家を賢くする}

\chapterintro{スマホには、ますます重要になる役割がもう一つあります——家じゅうの「リモコン」になることです。電球1個、スマートプラグ1つ、玄関の鍵にいたるまで、スマートホームは「ガジェット好きのおもちゃ」から「ごくふつうの設備」になりつつあります。この章では、プロトコル·エコシステム·ハブ·入門セットを整理し、最後に安全チェックリストで締めくくります。「スマート」に足をすくわれないための知識です。}

\booksection{スマートホームのプロトコル：どれを選ぶかが大事}

スマートデバイスを買う前に、まずどんなプロトコルに対応しているかを見ましょう。スマホやエコシステムがそのデバイスを「思いどおりに動かせる」かどうかが、これで決まります。いまの主流は4種類。それぞれの性格は次のとおりです：

\begin{itemize}
  \item Matter プロトコル：2022年に CSA（接続標準連盟）が発表。背後には Apple、Google、Amazon、Samsung などの大手が手を組んでいます。目的はただ一つ——すべてのブランドのスマートデバイスに「同じ言葉」を話させること。2023〜2026年には急速に統一標準になりつつあり、対応デバイスは Apple HomeKit、Google Home、Amazon Alexa、Mi Home（米家）に同時に接続でき、単一のエコシステムに縛られません。
  \item Thread プロトコル：Matter の土台となるネットワークの「配管」です。低消費電力でメッシュ構成（どこかのノードの電源が落ちても、信号が自動的に別の経路を通る）、応答が速く、もたつきません。Thread ボーダールーターを備えたデバイス（HomePod mini、Apple TV 4K、Xiaomi ハブゲートウェイなど）を選ぶと、動作が明らかに安定します。
  \item Wi-Fi：帯域が広く、ルーターに直結できますが、消費電力が高く、デバイスが増えるとチャンネルを圧迫しがち。カメラやロボット掃除機など、大きなデータをやり取りするデバイス向きです。
  \item Bluetooth（BLE）：初期設定が簡単で省電力ですが、距離が短く、メッシュ構成はできません。玄関鍵や温湿度計など、小さなデータ量のデバイス向きです。
  \item Zigbee：昔からある低消費電力のメッシュ型プロトコルです。専用ゲートウェイが必要なものの、デバイスは安価で安定。ただしブランドをまたぐ相互接続が弱く、Matter＋Thread に少しずつ置き換えられつつあります。
\end{itemize}

\begin{tcbinfo}{ひとことで言うと}
2026年にスマートデバイスを買うなら、まず「Matter over Thread」の表記があるかを確認——これが、いま最も時間の試練に耐える選択です。古いデバイスが Wi-Fi や Bluetooth までしか対応していなくても、急いで処分する必要はありません。使い続けて、新しく買うときに Matter へ向かえばOKです。

\end{tcbinfo}

\booksection{三大エコシステムの比較}

プロトコルは「言葉」で、エコシステムは「コミュニティ」です。どれも Matter に対応していても、日々のオートメーションや音声操作のためには、メインのエコシステムを1つ選ぶことになります。三大主流エコシステムを比較してみましょう：

\begin{center}
\begin{tabularx}{\linewidth}{|X|X|X|X|}
\hline
\textbf{エコシステム} & \textbf{強み} & \textbf{弱み} & \textbf{向いている人} \tabularnewline
\hline
Apple HomeKit & プライバシーが最強（ローカル処理が中心）、iPhone／Siri／ショートカットとの連携が滑らか & デバイスが少なめで高め、HomePod か Apple TV がハブとして必要 & 家じゅう Apple 製で、プライバシーと体験を重視するユーザー \tabularnewline
Mi Home（Xiaomi） & デバイスが最も豊富でコスパ最高、「Xiao AI（小愛同学）」も使いやすい & エコシステムがやや閉じ気味、データはクラウド処理が多く、広告も多め & 予算が限られていて、家じゅうを一度にそろえたい入門ユーザー \tabularnewline
HUAWEI スマートホーム（HarmonyOS） & HarmonyOS のスマホ／タブレット／ウォッチとシームレスに連携、対応デバイス数も急増中 & 海外デバイスの対応が少なく、エコシステムはまだ追い上げ中 & HUAWEI スマホのユーザー、「スーパーデバイス」による複数デバイス連携を重視する人 \tabularnewline
\hline
\end{tabularx}
\end{center}

\begin{tcbinfo}{エコシステム選びのアドバイス}
スマホに合わせてエコシステムを選ぶ：iPhone ユーザーは HomeKit を、HUAWEI ユーザーは HUAWEI スマートホームを、その他の Android ユーザーは Mi Home がいちばん楽です。Matter プロトコルのおかげで、別々のエコシステムのデバイスも同じ1台のスマホで操作できるので、最初から一つに「決め込む」必要はありません。

\end{tcbinfo}

\booksection{スマホをホームハブにする}

入門レベルの多くの場面では、スマホこそがあなたの「ホームハブ」です——アプリを通して照明·エアコン·カーテン·防犯を操作し、オートメーションを組み立てて、家を「自分で動くもの」にできます：

\begin{itemize}
  \item デバイス操作：HomeKit／Mi Home／HUAWEI AI Life アプリを開けば、遠隔でも自宅内でも、照明のオン／オフ、色温度の調整、エアコンの温度設定、カーテンの開閉、カメラ映像の確認ができます。
  \item オートメーション：「条件 動作」を組み合わせます。よく使う3大シーン： 帰宅モード：自宅への到着を検知＋日没 玄関灯を点ける＋エアコンを入れる＋カーテンを閉める＋玄関鍵を施錠。
  \item 外出モード：最後の1人が出かける（自宅を離れたことを検知、または玄関鍵の施錠） 照明を消す＋エアコンを切る＋スマートプラグを切る＋カメラを起動＋警戒モードに設定。
  \item おやすみモード：毎晩23:00 家じゅうの主照明を消す＋カーテンを閉める＋エアコンを26度に切り替え＋常夜灯を点ける。
\end{itemize}

NFC タッチ：玄関に NFC タグを貼っておけば、スマホをかざすだけで対応するシーンが動きます（1回タッチで玄関灯を点灯、もう1回でおやすみモードへ）。iPhone は「ショートカット」のオートメーションで、Android は「ウォレット」系アプリや各メーカーのアプリで書き込めます。

\begin{tcbpractice}{やってみよう：「帰宅モード」を組み立てる}
スマホのスマートホームアプリの「オートメーション」で新規作成します：トリガーは「特定の場所に到着」（自宅）を選び、アクションには順に「玄関灯を点ける」「エアコンを26度に設定」「カーテンを閉める」を追加。保存したら一度外へ出て、帰宅したときに勝手に明かりがつくか確かめましょう——初めて成功したときの達成感は、新しいデバイスをポチるより気持ちいいものです。

\end{tcbpractice}

\booksection{入門セットのおすすめ}

一度に全部そろえる必要はありません。この4点セットから始めれば、日常のニーズの8割はカバーできます：

\begin{itemize}
  \item スマートスピーカー：Xiao AI（Mi Home エコシステム、コスパが高い）／HomePod mini（Apple エコシステム、Thread ハブも兼ねる）／Alexa（海外の対応が広い）。注意：スピーカーで音声エコシステムが決まります。まずエコシステムを決めてからスピーカーを買いましょう。
  \item スマート電球：Yeelight（Mi Home／Matter 対応、色温度と明るさを調節できて安い）／Philips Hue（色がきれいで安定、ただし Hue ブリッジが必要でやや高価）。注意：買う前にソケットの規格（E27／B22）と、色調整の要否を確認しましょう。
  \item スマートプラグ：Xiaomi スマートプラグ（安価、電力計測つき）／TP-Link Kasa（海外向けの Matter 対応が良い）。注意：大電力の家電（給湯器、電気ヒーターなど）に使うときは、必ず定格電力を確認して過負荷を避けましょう。
  \item スマートロック：Lockin／Xiaomi（コスパ、Mi Home 連携）／Kaadas（錠前の伝統的な大手で、機械構造が堅牢）。注意：Matter や自分のエコシステムへの対応を確認し、電池のもちと緊急用キーホールの位置にも注目しましょう。
\end{itemize}

\begin{tcbinfo}{購入の順番のおすすめ}
まずスマートスピーカー（エコシステムを決める） 次にスマート電球とスマートプラグ（ハードルが低く、すぐ効果を実感できる） 最後に玄関鍵とカメラ（安全に関わるので、慎重に選ぶ）。一度に追加するのは1種類までにして、慣れてから増やせば、「買ったのに使えない」を避けられます。

\end{tcbinfo}

\booksection{セキュリティ上の注意点}

スマートデバイスが多い家ほど、攻撃の入り口も広がります。次の4つは最低ラインです。1つずつ照らし合わせてください：

\begin{enumerate}
  \item 初期パスワードを変える：カメラ、ルーター、玄関鍵の出荷時パスワードは、ネット上ですべて公開されています。買ったらまず、長くて使い回しのない強いパスワードに変更しましょう。
  \item ファームウェアをこまめに更新：メーカーが配るファームウェアには、セキュリティ修正が入っていることが多いもの。アプリで「自動更新」をオンにするか、少なくとも月1回は手動で確認しましょう。
  \item ネットワーク分離：IoT デバイスは専用の Wi-Fi（またはゲストネットワーク）に置き、スマホやパソコンのメインネットワークとは分けます。万が一どれかのデバイスが突破されても、大事なデータへの横展開の脅威を及ぼしにくくなります。
  \item プライバシー保護：カメラは「プライバシーゾーン」のマスキング（寝室や洗面所は撮らないなど）を有効にし、クラウドでの24時間録画はオフに。自宅カメラの映像を、サードパーティ製アプリに長期間見せっぱなしにしないこと。
\end{enumerate}

\textbf{この落とし穴にはハマらないで}\par

\begin{itemize}
  \item 知らない人にスマートロックの「遠隔解錠」を使わせない——一時的な暗証番号を優先し、使い終わったら失効させましょう。
  \item カメラをベッドや窓の外に直接向けない——漏れるのは自分のプライバシーです。
  \item 安さにつられて無名の「ノーブランド」デバイスを買わない——ファームウェアの更新がほとんどなく、家に裏口を取り付けるようなものです。
\end{itemize}

\begin{tcbsummary}{この章のまとめ}
\begin{itemize}
  \item Matter＋Thread が2023〜2026年の統一の流れ。新しく買うなら、まずこの表記をチェック
  \item 三大エコシステムにはそれぞれ得意分野がある：Apple はプライバシーと体験、Mi Home はコスパ、HUAWEI は複数デバイス連携
  \item スマホはハブになる。「帰宅／外出／おやすみ」の3大シーンを組み上げれば、家が自分で動き出します
  \item 入門の4点セット：スマートスピーカー、電球、スマートプラグ、玄関鍵。この順番で少しずつ追加
  \item 初期パスワードの変更、ファームウェア更新、ネットワーク分離、プライバシーのマスキング——4つの最低ラインは省けません
\end{itemize}

\end{tcbsummary}
